\section{Sampling integral columns}\label{sam:section}

We now combine the two residue constructions and lift them to integral
columns.  The purpose of this section is to express the probability of a
small determinant as a weighted lattice count.  The main and auxiliary
changes of coordinates are shared by all samples; fresh labels, digits,
auxiliary columns, and lifts are chosen for each sample.

Set
\begin{equation}\label{sam:box}
 N=(hq)^{10},\qquad
 \mathcal B=\bigl([0,N)^2\times[N,2N)\bigr)\cap\mathbb Z^3.
\end{equation}
For $u\in\mathcal B$, define its projection by
\[
 \pi(u)=(u_1/u_3,u_2/u_3)\in[0,1)^2.
\]
Choose $G_h$ uniformly in $\operatorname{SL}_3(\mathbb Z/h\mathbb Z)$,
independently of the auxiliary set $V$.  Conditional on $(G_h,V)$, sample
columns $u\in\mathcal B$ independently as follows:
\begin{enumerate}
 \item Draw a label uniformly from $K$ and its digit column $c$ by the
 main-modulus construction, using fresh independent digit choices.
 Prescribe $u\bmod h=G_hc$.
 \item Independently choose $u\bmod q$ uniformly from $V$.
 \item Among columns with those residues, choose $u$ uniformly in
 $\mathcal B$.
\end{enumerate}
Since $h$ and $q$ are coprime and $hq$ divides $N$, each pair of prescribed
residues has exactly $(N/(hq))^3$ lifts.  This also shows that the sampling
rule is defined for every outcome of the residue choices.  The samples
need not be independent after averaging over $(G_h,V)$; all estimates below
retain the shared choices.

\begin{lemma}[Proportional pairs]\label{sam:proportional}
For two sampled columns $u,v$, one has
\begin{equation}\label{sam:pair-bound}
 \Pr\bigl(\pi(u)=\pi(v)\bigr)\ll (hq)^6N^{-3}.
\end{equation}
\end{lemma}
\begin{proof}
Equality of the projections is equivalent to proportionality of $u$ and
$v$, because their third coordinates are positive.  Every column in
$\mathcal B$ is a positive integer multiple of a unique primitive integral
vector $z$ with $z_3>0$.  The number of its multiples in $\mathcal B$ is at
most $\sqrt6N/|z|$.  Therefore the total number of ordered proportional
pairs in $\mathcal B^2$, including equal columns, is at most
\[
 6N^2\sum_{0<|z|\le\sqrt6N}|z|^{-2}\ll N^3.
\]
For the last estimate, dyadic shells of radius $R$ contain $O(R^3)$
integral vectors and contribute $O(R)$ to the sum.

Condition on both columns' residues at both moduli and on the shared
random choices.  Their lifts are independent, each uniform on exactly
$(N/(hq))^3$ columns.  The preceding count bounds the number of
proportional pairs within any two such sets of lifts.  Dividing by
$(N/(hq))^6$ and averaging proves the claim.
\end{proof}

For three samples, write $A=(u^{(1)},u^{(2)},u^{(3)})$ for the integral
matrix of columns, and $C=(c^{(1)},c^{(2)},c^{(3)})$ for their digit matrix
before multiplication by $G_h$.  With $C$ fixed, use
Lemma~\ref{orb:diagonal} to define $D,E,I=DE$ and the row lattice
$\Lambda=\Lambda_C$, and write $\mathcal O=\mathcal O_C$ for the orbit in
Lemma~\ref{orb:orbit}.  Thus $\Lambda$ is the inverse image in $\mathbb Z^3$ of the row
span of $C$ modulo $h$, and
$\mathcal O=\{GC:G\in\operatorname{SL}_3(\mathbb Z/h\mathbb Z)\}$.
Recall also the auxiliary weight from \eqref{aux:weight}:
\[
 W(A)=\left(\frac{q^3}{s}\right)^3
 \Pr\bigl(u^{(1)},u^{(2)},u^{(3)}\bmod q\in V\bigr).
\]
Here the probability defining $W$ averages only over the auxiliary choices.

\begin{lemma}[Exact lifting identity]\label{sam:lifting}
Condition on all three labels and digit columns, hence on $C$.  For any
set $\mathcal E\subseteq\mathcal B^3$ of integral matrices,
\begin{equation}\label{sam:crt-identity}
 \Pr(A\in\mathcal E\mid C,\text{labels})
 =\frac{h^9}{N^9|\mathcal O|}
   \sum_{\substack{A\in\mathcal E\\ A\bmod h\in\mathcal O}}W(A).
\end{equation}
\end{lemma}
\begin{proof}
The matrix $G_hC$ is uniform on $\mathcal O$, since every point of an orbit
has the same number of group elements mapping $C$ to it.  Given $V$, the
three auxiliary residues are independent and uniform on its $s$ elements.
Consequently, for each specified ordered triple of auxiliary residues,
its probability is its inclusion probability in $V$, divided by $s^3$.
This remains true when some residues coincide.  Independence of the main
and auxiliary choices and the exact number of lifts give, for each fixed
integral $A$ with $A\bmod h\in\mathcal O$, the conditional probability
\[
 \frac{1}{|\mathcal O|}\,
 \frac{\Pr(A\bmod q\text{ has all columns in }V)}{s^3}
 \left(\frac{hq}{N}\right)^9
 =\frac{h^9}{N^9|\mathcal O|}W(A).
\]
For other matrices the probability is zero.  Summation proves the identity.
\end{proof}

We call a sampled triple \emph{bad} if its projected points are pairwise
distinct and $|\det A|\le\tau$.  Every matrix in $\mathcal O$ has determinant
$\det C$ modulo $h$.  Since $2\tau<h$, at most one integer
$t\in[-\tau,\tau]$ can occur as the determinant of a bad triple with
$A\bmod h\in\mathcal O$.  If the labels are all distinct,
Lemma~\ref{norm:obstruction} rules out every such $t$.

The following estimate handles the one remaining determinant value.
Its nonzero case follows immediately from the lattice count; the zero
case requires the auxiliary construction and is proved in the next
section.

\begin{proposition}[Weighted count at a fixed determinant]
\label{sam:weighted-count}
Fix a digit matrix $C$, with row lattice $\Lambda$ of index $I=DE$ and
special-linear orbit $\mathcal O$ modulo $h$.  For each integer
$t$ with $|t|\le\tau$, let $\mathcal A_t$ consist of the matrices $A$ whose
columns lie in $\mathcal B$, whose projected columns are pairwise
distinct, and which satisfy $A\bmod h\in\mathcal O$ and $\det A=t$.
Then
\begin{equation}\label{sam:weighted-estimate}
 \sum_{A\in\mathcal A_t}W(A)
 \ll (\log(2N))^2\frac{N^6}{I^2}.
\end{equation}
\end{proposition}
\begin{proof}
Suppose first that $t\ne0$.  Since $0<|t|\le\tau<q$, the reduction of
$A$ modulo $q$ has rank three.  Its columns are therefore affinely
independent, and Proposition~\ref{aux:weights} gives $W(A)\ll1$.
Every row of $A$ belongs to $\Lambda$ and has norm less than $4N$.
Moreover $E\mathbb Z^3\subseteq\Lambda$ gives
$\lambda_3(\Lambda)\le E\le h<4N$.  Proposition~\ref{lat:fixed-det},
applied with $X=4N$, proves the estimate.

For $t=0$, Proposition~\ref{zero:weighted-zero} below proves the stronger
bound $O(\log(2N)N^6/I^2)$ for all matrices with rows in $\Lambda$ and
pairwise distinct projected columns.  The orbit restriction only reduces
this set, so it gives the claimed estimate as well.
\end{proof}

\begin{corollary}[Probability of a bad triple]\label{sam:bad-triple}
Three sampled columns satisfy
\begin{equation}\label{sam:triple-bound}
 \Pr\bigl(\pi(u^{(1)}),\pi(u^{(2)}),\pi(u^{(3)})
       \text{ are pairwise distinct},\ |\det A|\le\tau\bigr)
 \ll (\log(2N))^2\frac{h}{N^3r^d}.
\end{equation}
\end{corollary}
\begin{proof}
Conditional on the labels and digit matrix $C$, the probability is zero
when all labels are distinct.  Otherwise Lemma~\ref{sam:lifting} and
Proposition~\ref{sam:weighted-count}, together with the bound
$|\mathcal O|\gg h^8/(D^3E^2)$ from \eqref{orb:orbit-bound}, bound it by
\[
 \frac{h^9}{N^9}\frac{D^3E^2}{h^8}
 (\log(2N))^2\frac{N^6}{D^2E^2}
 \ll (\log(2N))^2\frac{hD}{N^3}.
\]
Let $F$ be the event that two of the three labels agree.  The uniform
conditional moment estimate \eqref{orb:uniform-moment} gives
\[
 \mathbb E\bigl[D\mathbf1_F\bigr]
 =\mathbb E_{\mathrm{labels}}
   \bigl[\mathbf1_F\mathbb E(D\mid\mathrm{labels})\bigr]
 \ll\Pr(F)\le 3r^{-d}.
\]
This step uses the moment estimate for every fixed label triple, including
repeated labels; the corresponding digit draws remain independent.
Averaging the conditional bound proves the corollary.
\end{proof}
